\subsection{Model validation through targeted perturbations}

To validate the model, we reproduced perturbation experiments from the original \textit{Dendronotus iris} studies \cite{SakuraiKatz2016, sakurai2022bursting}.
Brief hyperpolarizing or depolarizing current pulses delivered during ongoing swim episodes test circuit resilience: a robust network should maintain coordinated output despite perturbations and recover rapidly upon their removal.
The model reproduces these experimental responses without parameter adjustment, confirming that it captures both the functional dynamics and stability properties of the biological CPG.

Figure~\ref{fig:perturbation_responses} presents six perturbation conditions; across all cases, the circuit recovers coordinated bursting once perturbations are removed, demonstrating that the network operates as a robust attractor.

Depolarizing Si3L produces immediate cessation of network bursting: the contralateral Si3R becomes hyperpolarized through reciprocal inhibition, and Si2L is silenced---losing excitatory drive from Si3R while receiving inhibition from the now-active Si2R.
Meanwhile, Si2R remains active, driven by the depolarized Si3L.
This perturbation reveals the asymmetric and distributed nature of control in the circuit.
Upon pulse termination, the circuit recovers rapidly.
Depolarizing Si3R similarly disrupts alternation; intermittent isolated spikes can be observed in the electrically coupled Si3L and Si2R neurons during the perturbation, indicating that electrical coupling allows small depolarizing excursions even under strong inhibitory pressure.
Upon pulse termination, the circuit rapidly recovers to its characteristic antiphase rhythm.

Hyperpolarizing a single Si2 neuron represents a special case: it is the only perturbation that does not completely halt network bursting.
Activity persists through alternative pathways---the contralateral E/I module maintains rhythmic drive despite suppression of its Si2 partner.
This response reveals a fundamental architectural principle: the contralateral E/I coupling constitutes the core oscillatory engine of the swim network, with half-center oscillators providing stabilization rather than rhythm generation.
Bilateral Si2 hyperpolarization permits recovery but with longer restoration time, highlighting the stabilizing role of Si2 reciprocal inhibition.
During bilateral Si2 suppression, small subthreshold oscillations persist in the inhibited neurons, likely reflecting residual drive through electrical coupling from the active Si3 neurons.

Hyperpolarizing Si3R silences its electrical partner Si2L through gap junction coupling.
Bilateral Si3 hyperpolarization abolishes network activity entirely: without Si3 excitatory drive, Si2 neurons---configured near the threshold between tonic spiking and quiescence---cannot sustain oscillation.
Upon release, the CPG rapidly restores its rhythm, confirming that Si3 neurons are necessary to initiate network-level bursting.
Additional bilateral perturbations further validate the model (Supplementary Materials).

These experiments assumed elevated Si3 excitability to sustain bursting.
In the biological circuit, this drive originates from Si1 command neurons, whose control over network state we examine next.
